\documentclass[10pt]{article}
\usepackage[a4paper,margin=1in]{geometry}
\usepackage{amsmath,amsthm,amssymb}
\usepackage{booktabs}
\usepackage{xcolor}
\definecolor{linkblue}{rgb}{0,0,0.5}
\usepackage[colorlinks=true,linkcolor=linkblue,citecolor=linkblue,urlcolor=linkblue,
  pdftitle={Explicit circuits for the Walsh-Hadamard transform below the textbook operation count},
  pdfauthor={Jacob Sussman}]{hyperref}

\newtheorem{theorem}{Theorem}
\renewcommand{\thetheorem}{\Alph{theorem}}
\newtheorem{proposition}{Proposition}
\newtheorem{lemma}[proposition]{Lemma}
\newtheorem{corollary}[proposition]{Corollary}
\theoremstyle{remark}
\newtheorem{remark}[proposition]{Remark}

\newcommand{\cntG}{\textup{(G)}}
\newcommand{\cntO}{\textup{(O)}}
\renewcommand{\lg}{\log_2}
\newcommand{\p}{$+$}
\newcommand{\m}{$-$}
\newcommand{\z}{$\cdot$}
\newcommand{\pc}[2]{(#1,\,#2)}
\hyphenation{Had-a-mard}
\newcommand{\reponame}{wht-fewer-operations}
\newcommand{\repourl}{https://github.com/jacobalansussman/\reponame}
\renewcommand{\labelitemi}{$\bullet$}
\emergencystretch=1.5em
\setlength{\textfloatsep}{12pt plus 2pt minus 2pt}
\setlength{\intextsep}{10pt plus 2pt minus 2pt}

\title{\vspace{-1.2cm}Explicit circuits for the Walsh--Hadamard transform\\ below the textbook operation count}
\author{Jacob Sussman}
\date{\vspace{-0.2cm}11 October 2026\\[3pt]{\normalsize Not peer reviewed; not submitted anywhere.}\vspace{-0.3cm}}

\begin{document}
\maketitle

\begin{abstract}
\noindent
The textbook program for the Walsh--Hadamard transform of $N=2^k$ numbers uses $N\lg N$ additions and
subtractions. I give explicit circuits that use fewer, in two ways of counting.
Count~\cntG: one operation is any $a u+b v$, and multiplying by a constant is free. Here the $8$-point
transform takes $19$ operations (textbook $24$; the smallest count I found in print is $22$), and a
recursive family takes $\tfrac34 N\lg N+\tfrac14 N$ for every $N\ge 4$: $52$ at $16$ points, $304$ at $64$.
The smallest constant in front of $N\lg N$ that I found in print is $7/8$ (Coppersmith, Feig and Linzer, 1994);
every line of my circuits is one operation in the count of that paper, so the comparison is like for like.
Count~\cntO: every addition, subtraction and multiplication by a constant costs one. Here the exact transform
on $N=8^j$ points takes at most $\tfrac56 N\lg N+N$ operations: $4{,}352$ against $4{,}608$ at $512$ points,
the first size where it is below the textbook ($28$ against $24$ at $8$ points).
The smallest constant that I found in print is $23/24$ (Alman and Rao, 2023).
The count-\cntO{} statement is checked in the proof assistant Lean. The others have short proofs, given here
in full; AI coding agents wrote them, and no person has checked them line by line. The circuits of all these
statements were replayed exactly by program up to $16{,}384$ points.
This is a better constant, not a better exponent: algorithms with $O(N(\log N)^{z})$ operations, $z<1$, have been
announced with proofs checked by Lean, and they are smaller for large~$N$. The circuits here are explicit and are small at sizes one can write
down. They were found by computer searches run with the same agents. ``New'' means that I did not find
them in the literature I searched. Not everything that bears on this was read:
Section~\ref{sec:found} says what was read in full, what in part and what from abstracts only.
A timing test on one machine found no speed-up; hardware figures are in Remark~\ref{rem:hardware}.
\end{abstract}

\section{Introduction}

\paragraph{The transform and the textbook count.}
\enlargethispage{\baselineskip}% one more line on the first page, for the last sentence of the abstract
Let $N=2^k$. Number the inputs $x_0,\dots,x_{N-1}$. The \emph{Walsh--Hadamard matrix} $H_N$ has, in row $i$ and
column $c$ (both numbered from $0$), the entry $+1$ if $i$ and $c$, written in binary, have an even number of
$1$-bits in common, and $-1$ otherwise. The \emph{Walsh--Hadamard transform} of $x$ is the list $H_N x$; I say
\emph{exact} when I mean $H_N x$ itself, with no factor in front. Since
$H_{2N}=\bigl(\begin{smallmatrix}H_N&H_N\\H_N&-H_N\end{smallmatrix}\bigr)$, the transform can be computed in $k$
layers, each of $N/2$ additions and $N/2$ subtractions. This is the textbook program, the fast
Walsh--Hadamard transform. It uses $N\lg N$ operations. When a count is $\gamma N\lg N$ plus a multiple
of $N$, I call $\gamma$ its \emph{rate}; the textbook has rate~$1$.

\paragraph{Two ways of counting.}
A \emph{circuit} is a list of lines. Each line computes one new value from inputs and from values of
earlier lines, and some of the values are named as the outputs. Every number in this paper belongs to one
of two counts.
\begin{itemize}\setlength{\itemsep}{1pt}\setlength{\parskip}{0pt}
\item \textbf{Count \cntG.} A line is $w=a u+b v$ with any constants $a,b$, and costs one. Multiplying by a
constant alone is free. This is the natural count for hardware in which a multiplication comes together
with the addition in one instruction (``multiply/add''), and for linear circuits whose gates have two
inputs, where it is half the number of wires (a \emph{wire} is one input of one gate).
\item \textbf{Count \cntO.} A line is $u+v$, $u-v$ or $c\,u$ with a constant $c$, and each costs one. This is
the number of arithmetic operations. It is the count of Alman and Rao~\cite{AR}, and the gate set of the
circuit model in OpenAI's Lean development~\cite{OAI1}.
\end{itemize}
The textbook costs $N\lg N$ in both counts. A circuit of count~\cntO{} is also a circuit of count~\cntG{} of the
same size, and one line of count~\cntG{} can be written with at most three lines of count~\cntO. No count-\cntO{}
circuit of this paper contains a bare change of sign, so it does not matter whether one charges for that.

\paragraph{The numbers.}
The inputs and the constants are rational, real or complex numbers. Proposition~\ref{prop:19} (its left
side) and Theorem~\ref{thm:B} use only the constants $1$, $-1$, $2$ and $-2$. Theorem~\ref{thm:A} and
Corollary~\ref{cor:allk} also use $\tfrac12$. Proposition~\ref{prop:lower} is stated for these three kinds of
numbers; its proof uses of them only that they form a field in which $2\ne0$.

\paragraph{What is in print.}
\emph{Count \cntG.} Coppersmith, Feig and Linzer~\cite{CFL} say on page~969 what they count: a multiply/add,
which ``computes $\pm a\pm bc$ where $a$, $b$ and $c$ are real numbers'', costs one, and so does a plain
addition or a plain multiplication. They give $\tfrac78\,n\lg n$ operations for the Hadamard transform of
length $n$ when $n$ is a power of four, and $n/8$ more when $n$ is an odd power of two. By these formulas
that is $7$ operations at $4$ points, $22$ at $8$ and $56$ at $16$. They print a $4$-point circuit and the
two formulas. They give no circuit of their own for $8$ or for $16$ points; for $16$ points they say that
eight $4$-point transforms can be used, with the same number of operations (page~970). They claim no lower
bound.

Count~\cntG{} allows a constant on both operands, which is more than a multiply/add. The circuits of this
paper do not use that: every line has two terms and at most one multiplication by a constant
(Remark~\ref{rem:cfl}). So each line is exactly one operation of~\cite{CFL}, and the comparisons with that
paper are like for like, in its own count: $19$ against $22$ at $8$ points, $52$ against $56$ at $16$, and
the rate $3/4$ against $7/8$.

Seven operations at $4$ points are due to Coppersmith, Feig and Linzer~\cite[page~969]{CFL}. They compute one
value that differs from three of the four outputs in a single sign, and get each of those outputs by adding
twice one input. A $4$-point circuit with $7$ operations was found again several times: by genetic
programming (Samadi, Suzukake and Iwakura~\cite{Samadi}, known from its abstract), in a patent of Reshef
and Alrod~\cite{Patent}, and by Alman and Rao~\cite{AR}, who use $7$ additions or subtractions and one
halving. Feig wrote, in a chapter dated 1999, that the $4$-point result ``has not been extended, and it may in fact be rather
hard to extend except in most trivial (and uninteresting) ways''~\cite{Feig}. The smallest count at $8$ points
that I found in print is $22$, and it appears twice. The formula of~\cite{CFL} for odd powers of two gives
$22$ at $8$ points: one textbook layer of $8$ operations and two $4$-point circuits of $7$ (they print the
formula, not the number). And the $8$-point step of Alman and Rao uses $22$ additions or subtractions and
one halving. The halving is free in count~\cntG, so their $23$ operations of count~\cntO{} are $22$ of count~\cntG.

\emph{Count \cntO.} Alman and Rao~\cite[Theorem~1.1]{AR} prove that for any field and any $N$ that is a power
of two, $H_N x$ can be computed with at most $\tfrac{23}{24}N\lg N+\tfrac{13}{12}N$ field operations. They
write that, to their knowledge, this is the first improvement on the operation count $N\log N$ of the
textbook program. The proof of their Theorem~6.2 gives the exact count
$\tfrac{23}{24}N\lg N+\tfrac{N}{24}(\lg N \bmod 3)+N-1$. Their method writes $H_8$ as the sum of a matrix
of low rank (one whose rows are all combinations of a few rows) and a sparse matrix (one with few non-zero
entries), and halves once; the resulting step costs $23$ operations per $8$ points. These
statements are in version~2 of their arXiv text. The idea of a single halving, used below, is theirs.

\emph{Other measures.} Alman~\cite{Alman} and later papers count the wires of circuits that are built in
layers and whose gates may have many inputs; a value that passes a layer unchanged costs a wire there. His
abstract gives a bound for every fixed number of layers and, with no limit on their number,
$(1.81+o(1))\,N\lg N$ wires. That is a third count, which I do not treat, and the numbers of this paper do
not convert into it line by line; Alman and Rao write
that the algorithms of~\cite{Alman} ``do not give improved arithmetic operation counts''~\cite{AR}.

\emph{The count of Grigoryan and Agaian.}
Grigoryan and Agaian~\cite{GA} state that their algorithm computes the Hadamard transform on $N=2^r$ points
with $6\cdot2^r-r^2-3r-6$ additions and subtractions (page~181), ``no more than six additions per sample''
(page~179). That count is too small. As the program printed on page~181 carries it out, the algorithm
applies a ``paired transform'' of length $N$, which costs $2N-2$ additions and subtractions, and this splits
the transform into transforms of lengths $N/2,N/4,\dots,2,1$, which are computed the same way. So the number
$T(N)$ of its operations obeys
\[
T(N)=(2N-2)+T(N/2)+T(N/4)+\dots+T(1),\qquad T(1)=0 ,
\]
and the solution is $T(N)=N\lg N$ exactly (by induction, because $\sum_{i<r}i\,2^i=(r-2)\,2^r+2$). The paper
does not print this recursion; it is the count of the algorithm, written down from that page. The stated
count agrees with it at $2$, $4$ and $8$ points ($2$, $8$, $24$) and is smaller from $16$ points on ($62$
against $64$). From $4{,}096$ points on it is also below $\tfrac12N\lg N$ ($24{,}390$ against $24{,}576$).
That number is a lower bound for circuits that use only additions and subtractions, by the determinant
argument of Morgenstern~\cite{Morgenstern}, taken here as Alman and Rao quote it~\cite{AR} (the
determinant of $H_N$ is $\pm N^{N/2}$); and the introduction of~\cite{GA} itself quotes $0.5\,N\log N$ as a
lower bound for the Hadamard transform (page~172). So the algorithm has the textbook count in another order,
and the count stated in~\cite{GA} is not a count below $N\lg N$. The program \verb|check.py| of
Section~\ref{sec:checked} prints the stated count beside the two other numbers for every size up to
$16{,}384$ points. A program that the agents wrote earlier from a later description of the algorithm by the
first author~\cite{GB} also used exactly $N\lg N$ additions and subtractions at every size from $2$ to
$4{,}096$ points; that program is not in the repository. Frigo and Johnson~\cite{FJ} published a comment on
the part of~\cite{GA} that concerns the Fourier transform. Only its abstract was read, and it does not speak
of the count for the Hadamard transform. Nothing is said here about the counts of~\cite{GA} for the
Fourier transform.

\emph{The exponent.} Two preprints of OpenAI of 25 September 2026~\cite{OAI1,OAI2} announce, with a Lean development, that the exact
discrete Fourier transform of length $n$ can be computed in $O(n(\log n)^{z})$ operations with a fixed $z<1$.
A proof, checked by Lean, for the Walsh--Hadamard and Fourier transforms with
$1-z\approx 8.7625\times 10^{-4}$ is in an earlier public repository of mine~\cite{RepoPS}, which this paper
does not use. Both count the steps of a machine
with exact arithmetic, not the gates of a circuit, so I compare orders of growth only. For large $N$ these are smaller
than $c\,N\log N$ for every constant $c>0$, and so smaller than every count in this paper.

\paragraph{The results.}
Table~\ref{tab:results} lists them. In this paper ``new'' always means: I did not find it in the literature
I searched. Section~\ref{sec:found} says what was searched, and what bears on the question and was not read.

\begin{table}[ht]
\centering
\small
\begin{tabular}{@{}lllll@{}}
\toprule
statement & count, sizes & this paper & found in print & how it is checked\\
\midrule
Prop.~\ref{prop:19} & \cntG, $N=8$ & $19$ & $22$ & replay; Table~\ref{tab:h8}\\
Thm.~\ref{thm:B} & \cntG, every $N=2^k\ge 4$ & $\tfrac34N\lg N+\tfrac14N$ & $\le\tfrac78N\lg N+\tfrac18N$ \cite{CFL}
   & proof; replay\\
Thm.~\ref{thm:A} & \cntO, $N=8^j$ & $\tfrac56N\lg N+N$ & $\tfrac{23}{24}N\lg N+N-1$ \cite{AR}
   & proof; Lean; replay\\
Cor.~\ref{cor:allk} & \cntO, every $N=2^k\ge 8$ & $\le\tfrac56N\lg N+\tfrac43N$ & $\le\tfrac{23}{24}N\lg N+\tfrac{13}{12}N$ \cite{AR}
   & proof; replay\\
Prop.~\ref{prop:lower} & \cntG{} and \cntO, $N\ge 4$ & at least $2N-1$ & not searched & proof\\
\bottomrule
\end{tabular}
\caption{The results, all for the exact transform $H_Nx$. ``Found in print'' is the smallest count
that I found in the literature I searched; for Theorem~\ref{thm:B} it is $\tfrac78N\lg N$ when $k$ is even
and $\tfrac18N$ more when $k$ is odd. ``Replay'' means: a program generated the circuit and ran it in
exact arithmetic on every \emph{unit input} (one input $1$, the others $0$), or made one run that is
equivalent to this, for every size up to $N=2^{14}$. Proposition~\ref{prop:19} is short enough to check by
hand from Table~\ref{tab:h8}. Section~\ref{sec:checked} says who wrote and who read the proofs.}
\label{tab:results}
\end{table}

In count~\cntO{} the circuits are below the textbook from $512$ points on. In count~\cntG{} they are below it from
$4$ points on, and below the counts I found in print from $8$ points on.

The circuits as text files, the programs that check them, the Lean files and the source of this paper are in
one repository~\cite{Repo}. ``The repository'' always means that one.

\section{The eight-point circuit}\label{sec:eight}

Table~\ref{tab:h8} gives the circuit. The middle columns show each value as its eight coefficients in
$x_0,\dots,x_7$.

\begin{table}[ht]
\centering
\small
\setlength{\tabcolsep}{4.2pt}
\begin{tabular}{@{}l@{\qquad}cccccccc@{\qquad}l@{}}
\toprule
line, count \cntG & $x_0$ & $x_1$ & $x_2$ & $x_3$ & $x_4$ & $x_5$ & $x_6$ & $x_7$ & centre halved, count \cntO\\
\midrule
$g_{1}=x_{0}+x_{7}$ & \p&\z&\z&\z&\z&\z&\z&\p & $g_{1}=x_{0}+x_{7}$\\
$g_{2}=x_{1}+x_{6}$ & \z&\p&\z&\z&\z&\z&\p&\z & $g_{2}=x_{1}+x_{6}$\\
$g_{3}=x_{2}-x_{5}$ & \z&\z&\p&\z&\z&\m&\z&\z & $g_{3}=x_{2}-x_{5}$\\
$g_{4}=x_{3}-x_{4}$ & \z&\z&\z&\p&\m&\z&\z&\z & $g_{4}=x_{3}-x_{4}$\\
$g_{5}=g_{1}-g_{2}$ & \p&\m&\z&\z&\z&\z&\m&\p & $g_{5}=g_{1}-g_{2}$\\
$g_{6}=g_{3}+g_{4}$ & \z&\z&\p&\p&\m&\m&\z&\z & $g_{6}=g_{3}+g_{4}$\\
$g_{7}=g_{5}-g_{6}$ & \p&\m&\m&\m&\p&\p&\m&\p & $g_{7}=g_{5}-g_{6}$, then $h=\tfrac12 g_{7}$\\
\midrule
$g_{8}=g_{7}-2x_{7}$ & \p&\m&\m&\m&\p&\p&\m&\m & $h_{8}=h-x_{7}$\\
$g_{9}=g_{7}-2x_{5}$ & \p&\m&\m&\m&\p&\m&\m&\p & $h_{9}=h-x_{5}$\\
$g_{10}=g_{8}+2x_{6}$ & \p&\m&\m&\m&\p&\p&\p&\m & $h_{10}=h_{8}+x_{6}$\\
$g_{11}=g_{9}-2x_{4}$ & \p&\m&\m&\m&\m&\m&\m&\p & $h_{11}=h_{9}-x_{4}$\\
$y_{7}=g_{10}+2g_{4}$ & \p&\m&\m&\p&\m&\p&\p&\m & $z_{7}=h_{10}+g_{4}$\\
$y_{1}=g_{10}+2g_{3}$ & \p&\m&\p&\m&\p&\m&\p&\m & $z_{1}=h_{10}+g_{3}$\\
$y_{6}=g_{11}+2g_{2}$ & \p&\p&\m&\m&\m&\m&\p&\p & $z_{6}=h_{11}+g_{2}$\\
$y_{0}=-g_{11}+2g_{1}$ & \p&\p&\p&\p&\p&\p&\p&\p & $z_{0}=g_{1}-h_{11}$\\
$y_{2}=g_{8}+2x_{1}$ & \p&\p&\m&\m&\p&\p&\m&\m & $z_{2}=h_{8}+x_{1}$\\
$y_{3}=g_{9}+2x_{3}$ & \p&\m&\m&\p&\p&\m&\m&\p & $z_{3}=h_{9}+x_{3}$\\
$y_{4}=y_{2}+2g_{6}$ & \p&\p&\p&\p&\m&\m&\m&\m & $z_{4}=z_{2}+g_{6}$\\
$y_{5}=-y_{3}+2g_{5}$ & \p&\m&\p&\m&\m&\p&\m&\p & $z_{5}=g_{5}-z_{3}$\\
\bottomrule
\end{tabular}
\caption{The $8$-point transform in $19$ operations of count~\cntG{} (left), and $\tfrac12H_8x$ in $20$
operations of count~\cntO{} (right). Middle: the coefficients of the value on the left; \p{} is $+1$, \m{} is
$-1$, \z{} is $0$. On the right the last twelve values are half of those on the left: $h_i=\tfrac12g_i$
and $z_i=\tfrac12y_i$.}
\label{tab:h8}
\end{table}

\begin{proposition}\label{prop:19}
The $19$ lines on the left of Table~\ref{tab:h8} compute the exact transform: $y_i$ is row $i$ of $H_8$
applied to $x$, for $i=0,\dots,7$. The $20$ lines on the right ($19$ additions or subtractions and one
halving) compute $z_i=\tfrac12 y_i$.
\end{proposition}

\begin{proof}
Check the table row by row. A line has two terms: its first operand, with a minus sign in two lines, and
its second operand times the constant shown. The pattern of the line is the sum of the patterns of the two
terms. In each of the last twelve lines the second term is non-zero only in columns where the first term has
the opposite sign, and there it is $\pm2$; so the line is the first term with exactly those signs flipped.
The eight rows named $y_i$ are the rows of $H_8$: for example $y_5$ has $-1$ exactly
in the columns $1,3,4,6$, the numbers that share an odd number of $1$-bits with $5=101_2$. On the right, the
first seven lines are unchanged, $h=\tfrac12g_7$, and each later line is the line on its left divided by~$2$.
\end{proof}

\paragraph{What the structure is.}
Seven lines merge the eight inputs into one value, $g_7$, which I call the \emph{centre}. It reads all
eight inputs and it is not an output. Each of the other twelve lines flips signs, by adding twice an input
or twice a value that the merging left behind ($g_1,\dots,g_6$). So every partner in those twelve lines was
already paid for. In the end every output is the centre, with a sign, plus twice a small sum:
\[
\begin{aligned}
y_2&=g_7+2(x_1-x_7), & y_0&=g_7+2(x_1+x_2+x_3+x_6),\\
y_3&=g_7+2(x_3-x_5), & y_1&=g_7+2(x_2-x_5+x_6-x_7),\\
y_4&=-g_7+2(x_0-x_6), & y_6&=g_7+2(x_1-x_4-x_5+x_6),\\
y_5&=g_7+2(x_2-x_4), & y_7&=g_7+2(x_3-x_4+x_6-x_7).
\end{aligned}
\]
Two of the $19$ lines, $y_0$ and $y_5$, have a minus sign on their first operand. All constants are $\pm1$ and
$\pm2$, and all $19$ values have coefficients in $\{0,+1,-1\}$.

In count~\cntO{} the $19$ lines as they stand would cost $31$: $19$ additions or subtractions and $12$
doublings. Halving the centre once, right after the line $g_7$, removes every doubling. This single halving
is the device of Alman and Rao~\cite{AR}; their $8$-point step costs $23$ operations where this one costs $20$.

\section{Count \cntO: the exact transform in \texorpdfstring{$\tfrac56N\lg N+N$}{(5/6) N log N + N} operations}\label{sec:A}

\begin{theorem}\label{thm:A}
Let $j\ge 1$ and $N=8^j$. There is a circuit with at most
\[
20\,j\,8^{\,j-1}+8^{\,j}\;=\;\tfrac56\,N\lg N+N
\]
operations of count~\cntO{} that computes the exact transform $H_N x$.
\end{theorem}

\begin{proof}
Write the number of a point with $3j$ binary digits, and cut them into $j$ groups of three. Row $r$,
column $c$ of $H_N$ is $(-1)^{\text{number of common $1$-bits}}$, and the common $1$-bits are counted group
by group. So $H_N$ is the product of $j$ matrices $L_1,\dots,L_j$ that commute, where $L_i$ acts as $H_8$ on
the $i$-th group of digits and leaves the other digits alone: $L_i$ splits the $N$ points into $N/8$ sets of
eight points that differ only in the $i$-th group, and applies $H_8$ to each set. (In matrix terms,
$H_N=H_8\otimes\dots\otimes H_8$ and $L_i=I\otimes\dots\otimes H_8\otimes\dots\otimes I$.)

Now apply the $20$-operation circuit of Proposition~\ref{prop:19} in place of $H_8$: for $i=1,\dots,j$ in turn,
run it on each of the $N/8$ sets of the $i$-th group, with the current values as inputs. This is $j$ layers of
$20\cdot N/8$ operations, and by Proposition~\ref{prop:19} it computes
$\tfrac12L_j\cdots\tfrac12L_1x=2^{-j}H_Nx$. Multiply each of the $N$ values by $2^{j}$. The total is
$20\,j\,N/8+N=20\,j\,8^{j-1}+8^{j}$ operations, and $20\,j\,N/8=\tfrac{20}{24}N\cdot 3j=\tfrac56N\lg N$.
\end{proof}

\begin{corollary}\label{cor:allk}
Let $k=3m+r$ with $m\ge1$ and $r\in\{0,1,2\}$, and $N=2^k$. Then $H_Nx$ can be computed exactly with
$N\,(\tfrac52m+r+1)$ operations of count~\cntO. This is $\tfrac56N\lg N+N$ for $r=0$,
$\tfrac56N\lg N+\tfrac76N$ for $r=1$ and $\tfrac56N\lg N+\tfrac43N$ for $r=2$. It is below the textbook's
$N\lg N$ exactly when $m\ge 3$, that is, from $512$ points on.
\end{corollary}

\begin{proof}
As in the proof of Theorem~\ref{thm:A}, use $m$ layers of the $20$-operation circuit on $3m$ of the digits
($\tfrac{20}{8}mN$ operations), one textbook layer for each of the $r$ other digits ($rN$ operations), and $N$
multiplications by $2^m$. Since $m=(k-r)/3$, the total is $\tfrac56Nk+(1+\tfrac r6)N$. It is smaller than
$Nk=N(3m+r)$ exactly when $\tfrac52m+1<3m$.
\end{proof}

\begin{table}[ht]
\centering
\small
\begin{tabular}{@{}rrrrr@{}}
\toprule
$N$ & textbook $N\lg N$ & Alman and Rao, exact count & this paper (Cor.~\ref{cor:allk}) & ratio to textbook\\
\midrule
8 & 24 & 30 & 28 & 1.167\\
16 & 64 & 77 & 72 & 1.125\\
32 & 160 & 187 & 176 & 1.100\\
64 & 384 & 431 & 384 & 1.000\\
128 & 896 & 991 & 896 & 1.000\\
256 & 2,048 & 2,239 & 2,048 & 1.000\\
512 & 4,608 & 4,927 & \textbf{4,352} & 0.944\\
1,024 & 10,240 & 10,879 & \textbf{9,728} & 0.950\\
2,048 & 22,528 & 23,807 & \textbf{21,504} & 0.955\\
4,096 & 49,152 & 51,199 & \textbf{45,056} & 0.917\\
\bottomrule
\end{tabular}
\caption{Count~\cntO, exact $H_Nx$. ``Alman and Rao, exact count'' is
$\tfrac{23}{24}N\lg N+\tfrac N{24}(\lg N\bmod 3)+N-1$, from the proof of their Theorem~6.2. Bold: below the
textbook.}
\label{tab:O}
\end{table}

\paragraph{The comparison.}
Table~\ref{tab:O} has the numbers from $8$ to $4{,}096$ points. Three things can be read from the formulas.
\begin{itemize}\setlength{\itemsep}{1pt}\setlength{\parskip}{0pt}
\item The count of Corollary~\ref{cor:allk} is below the textbook at every $N\ge 512$, equal to it at
$64$, $128$ and $256$ points, and above it at $8$, $16$ and $32$ points. Below $512$ points the textbook is
the smallest exact count in count~\cntO{} that I know.
\item It is below Alman and Rao's exact count at every $N\ge 8$: the difference is $\tfrac38mN-1$.
\item Alman and Rao's exact count is above the textbook at every size in the table. By its formula it is
first below the textbook at $N=2^{24}$, by one operation; their stated bound
$\tfrac{23}{24}N\lg N+\tfrac{13}{12}N$ is first below it at $N=2^{27}$. Their theorem is stated for any field.
Their circuit and mine both divide by $2$, so both need $2$ to be different from $0$ (where $2=0$ every entry
of $H_N$ is $1$ and the transform is the sum of the inputs). The Lean statement is over the complex numbers.
\end{itemize}

\begin{remark}[if a common factor is allowed]\label{rem:factor}
Without the last $N$ multiplications the same circuits compute $2^{-m}H_Nx$ in $N(\tfrac52m+r)$ operations of
count~\cntO: $20$ at $8$ points, $56$ at $16$, $320$ at $64$, $3{,}840$ at $512$. This is below $N\lg N$ from
$8$ points on. For even $k$ the family of Section~\ref{sec:B}, with its pieces kept at suitable powers of
$\tfrac12$, gives $H_Nx/\sqrt N$ in $\tfrac78N\lg N+\tfrac14N$ operations of count~\cntO, $60$ at $16$ points; this was
replayed by the programs of the repository~\cite{Repo} up to $16{,}384$ points, and I do not prove it here.

That rate, $7/8$, is above $5/6$. On 11 October 2026 the agents of Section~\ref{sec:found} looked for a way
to bring the rate $3/4$ of the family into count~\cntO, and found none. First, the lines of the family were
kept and each value was held multiplied by a power of $2$ that could be chosen freely; a value that is
needed with two different powers costs one more operation. The least count of this kind, with the exponents
taken from a fixed range, was computed exactly at $4$ to $64$ points. It is the $60$ above at $16$ points and $352$ at $64$, against $56$ and $320$
for the layers of this section, and an argument puts it at $\tfrac78N\lg N+\tfrac18N$ or more at every size.
Second, six ways of putting levels of the family together with the $8$-point circuit were built and
replayed, up to $4{,}096$ points; none costs less than the layers. These are computations and an argument of
the agents, not in Lean and not checked by a person, and they say nothing about circuits that are wired in another
way. In particular, no direct search has been made for a $16$-point circuit that computes a multiple of
$H_{16}x$ with $53$ or fewer operations of count~\cntO; used in layers, such a circuit would have a rate
below $5/6$. The best I know has $56$.
\end{remark}

\section{Count \cntG: a family with \texorpdfstring{$\tfrac34N\lg N+\tfrac14N$}{(3/4) N log N + N/4} operations}\label{sec:B}

\begin{theorem}\label{thm:B}
For every $N=2^k$ with $k\ge2$ there is a circuit with exactly
\[
\tfrac34\,N\lg N+\tfrac14\,N
\]
operations of count~\cntG{} that computes the exact transform $H_Nx$. All its constants are $\pm1$ and~$\pm2$.
\end{theorem}

The counts are $7$, $20$, $52$, $128$ and $304$ for $N=4,\dots,64$; Table~\ref{tab:G} has more. At $4$ points the circuit has the $7$ operations of Coppersmith, Feig and Linzer~\cite{CFL}, in another
wiring of the same idea: theirs has $4$ additions or subtractions and $3$ lines with a doubling, and the
family at $4$ points has $3$ and~$4$.
At $8$ points the family has $20$ operations, one more than the circuit of Section~\ref{sec:eight}; from
$16$ points on it is the smallest count~\cntG{} that I know.

\paragraph{Words.}
A \emph{block} of size $m=2^i$ is a run of $m$ consecutive inputs $x_e,\dots,x_{e+m-1}$ that starts at a
multiple $e$ of $m$. For a block $X$ of size $m$ and $0\le a<m$, the \emph{piece} $\pc Xa$ is row $a$ of the
$m$-point transform applied to the block:
\[
\pc Xa\;=\;\sum_{t=0}^{m-1}(-1)^{\text{number of common $1$-bits of $a$ and $t$}}\;x_{e+t}.
\]
A block of size $1$ has one piece, the input itself. The $N$ outputs are the $N$ pieces of the block of
size~$N$. A block $X$ of size $m\ge2$ has a lower half $L$ and an upper half $R$, and for $0\le a<m/2$
\begin{equation}\label{eq:halves}
\pc Xa=\pc La+\pc Ra,\qquad \pc X{a+m/2}=\pc La-\pc Ra .
\end{equation}
This is the textbook step. The textbook makes all $m$ pieces of every block, and so keeps every piece of
$L$ and every piece of $R$. The family keeps half of them.

\begin{lemma}[the hub]\label{lem:hub}
Let $X$ be a block of size $m\ge4$ with halves $L,R$ and quarters $L_0,L_1,R_0,R_1$, in this order. Let
$0\le d<m/4$, and write $q(Y)=\pc Yd$ for a quarter $Y$. Put
\[
A=\pc Ld,\qquad B=\pc R{d+m/4},\qquad M=A+c\,B\quad\text{with } c\in\{+1,-1\}.
\]
Let $s\in\{+1,-1\}$, and let $P_1=\pc Ld+s\,\pc Rd$ and $P_2=\pc L{d+m/4}+s\,\pc R{d+m/4}$. Then
\begin{align}
P_1&=M+(s-c)\,q(R_0)+(s+c)\,q(R_1), \label{eq:P1}\\
P_2&=cs\,M+(1-cs)\,q(L_0)-(1+cs)\,q(L_1). \label{eq:P2}
\end{align}
So if $c=-s$ then $P_1=M+2s\,q(R_0)$ and $P_2=-M+2\,q(L_0)$, and if $c=s$ then $P_1=M+2s\,q(R_1)$ and
$P_2=M-2\,q(L_1)$. By \eqref{eq:halves}, $P_1$ and $P_2$ are the pieces $\pc Xd$ and $\pc X{d+m/4}$ of $X$ when
$s=+1$, and the pieces $\pc X{d+m/2}$ and $\pc X{d+m/4+m/2}$ when $s=-1$.
\end{lemma}

\begin{proof}
By \eqref{eq:halves} applied to $L$ and to $R$, whose halves are the quarters,
\[
\begin{aligned}
\pc Ld&=q(L_0)+q(L_1), & \pc L{d+m/4}&=q(L_0)-q(L_1),\\
\pc Rd&=q(R_0)+q(R_1), & \pc R{d+m/4}&=q(R_0)-q(R_1).
\end{aligned}
\]
Hence $M=q(L_0)+q(L_1)+c\,q(R_0)-c\,q(R_1)$, while $P_1=q(L_0)+q(L_1)+s\,q(R_0)+s\,q(R_1)$ and
$P_2=q(L_0)-q(L_1)+s\,q(R_0)-s\,q(R_1)$. Subtracting $M$ from $P_1$ gives \eqref{eq:P1}. Since $c^2=1$,
$cs\,M=cs\,q(L_0)+cs\,q(L_1)+s\,q(R_0)-s\,q(R_1)$, and subtracting it from $P_2$ gives \eqref{eq:P2}.
\end{proof}

I call $M$ a \emph{hub}. It is not a piece: it is one piece of $L$ on the lower half plus a different piece
of $R$ on the upper half. It is one operation away from two pieces of $X$, and the partners in those two
operations are pieces of quarters.

\paragraph{The procedure $E(X,s)$.}
For a block $X$ of size $m\ge2$ and a sign $s$, the procedure $E(X,s)$ makes the $m/2$ pieces
$\pc La+s\,\pc Ra$, $0\le a<m/2$. By \eqref{eq:halves} these are the pieces $\pc Xa$ with $a<m/2$ if $s=+1$,
and with $a\ge m/2$ if $s=-1$.
\begin{itemize}\setlength{\itemsep}{1pt}\setlength{\parskip}{0pt}
\item If $m=2$: one line, $x_e+s\,x_{e+1}$.
\item If $m\ge4$: first run $E(L,+1)$ and $E(R,-1)$. Then, for every $d$ with $0\le d<m/4$, write three lines
in the notation of Lemma~\ref{lem:hub}: the hub $M=A+cB$ with $c=-s$ if $d<m/8$ and $c=s$ if $d\ge m/8$;
then $P_1$; then $P_2$, each from $M$ and one quarter piece as in the lemma.
\end{itemize}

\begin{lemma}\label{lem:E}
For every block $X$ of size $m=2^i\ge2$ and every sign $s$, each line of $E(X,s)$ uses only inputs and
values of earlier lines, $E(X,s)$ makes the $m/2$ pieces stated, and it has exactly $E(m)=m(3i-1)/4$ lines.
\end{lemma}

\begin{proof}
By induction on $m$. For $m=2$ the line is $\pc L0+s\,\pc R0$, and $E(2)=1$. Let $m\ge4$ and suppose the
lemma holds for $m/2$ and, if $m\ge 8$, for $m/4$. Fix $d<m/4$. The piece $A=\pc Ld$ has index
$d<\tfrac12(m/2)$, so $E(L,+1)$ made it. The piece $B=\pc R{d+m/4}$ has index at least $\tfrac12(m/2)$, so
$E(R,-1)$ made it. For the quarter pieces there are two cases. If $m=4$, then $d=0<m/8$, the quarters are
single inputs, and $q(L_0)$, $q(R_0)$ are inputs. If $m\ge8$, then $E(L,+1)$ began by running $E(L_0,+1)$
and $E(L_1,-1)$, and $E(R,-1)$ began by running $E(R_0,+1)$ and $E(R_1,-1)$. By induction these made
$\pc{L_0}d$ and $\pc{R_0}d$ for $d<m/8$, and $\pc{L_1}d$ and $\pc{R_1}d$ for $m/8\le d<m/4$. So for $d<m/8$ the
pieces $q(L_0),q(R_0)$ are at hand and, with $c=-s$, Lemma~\ref{lem:hub} gives $P_1=M+2s\,q(R_0)$ and
$P_2=-M+2\,q(L_0)$. For $d\ge m/8$ the pieces $q(L_1),q(R_1)$ are at hand and, with $c=s$, it gives
$P_1=M+2s\,q(R_1)$ and $P_2=M-2\,q(L_1)$. Each is one line of count~\cntG. As $d$ runs from $0$ to $m/4-1$, the
pieces $P_1,P_2$ are $\pc La+s\,\pc Ra$ for $a=d$ and $a=d+m/4$, which is every $a<m/2$ once.
The number of lines is $2E(m/2)+3\cdot\tfrac m4=\tfrac m4\bigl(3(i-1)-1\bigr)+\tfrac{3m}4=\tfrac m4(3i-1)$.
\end{proof}

\begin{proof}[Proof of Theorem~\ref{thm:B}]
Let $X$ be the block of all $N$ inputs, with halves $L,R$. Run $E(L,+1)$ and $E(R,-1)$. For every $d$ with
$0\le d<N/4$ write five lines, in the notation of Lemma~\ref{lem:hub} with $m=N$ and $c=+1$, so that $M=A+B$.
\begin{itemize}\setlength{\itemsep}{1pt}\setlength{\parskip}{0pt}
\item If $d<N/8$: $U=M-2\,q(R_0)$ and $V=-M+2\,q(L_0)$; then $2A-U$ and $V+2B$.
\item If $d\ge N/8$: $U'=M+2\,q(R_1)$ and $V'=M-2\,q(L_1)$; then $2A-U'$ and $V'-2B$.
\end{itemize}
The values $A$, $B$ and the quarter pieces are at hand for the same reasons as in the proof of
Lemma~\ref{lem:E}. In the first case Lemma~\ref{lem:hub} with $s=-1=-c$ gives $U=\pc Ld-\pc Rd=\pc X{d+N/2}$
and $V=\pc L{d+N/4}-\pc R{d+N/4}=\pc X{d+N/4+N/2}$. Then $2A-U=\pc Ld+\pc Rd=\pc Xd$, and
$V+2B=\pc L{d+N/4}+\pc R{d+N/4}=\pc X{d+N/4}$, because $B=\pc R{d+N/4}$. In the second case the lemma with
$s=+1=c$ gives $U'=\pc Xd$ and $V'=\pc X{d+N/4}$, and then $2A-U'=\pc X{d+N/2}$ and $V'-2B=\pc X{d+N/4+N/2}$.
So the four outputs with the numbers $d$, $d+N/4$, $d+N/2$, $d+3N/4$ are made, and as $d$ runs from $0$ to
$N/4-1$ these are all $N$ outputs. Every constant in these lines is $\pm1$ or $\pm2$. By Lemma~\ref{lem:E} the
number of lines is
\[
2E(N/2)+5\cdot\tfrac N4=\tfrac N4\bigl(3(k-1)-1\bigr)+\tfrac{5N}4=\tfrac N4(3k+1)=\tfrac34N\lg N+\tfrac14N .
\qedhere
\]
\end{proof}

\begin{table}[!htb]
\centering
\small
\begin{tabular}{@{}rrrrr@{}}
\toprule
$N$ & textbook & rate $7/8$ found in print & layers of the $19$ & family (Theorem~\ref{thm:B})\\
\midrule
4 & 8 & 7 & 7 & 7\\
8 & 24 & 22 & \textbf{19} & 20\\
16 & 64 & 56 & 54 & \textbf{52}\\
32 & 160 & 144 & 132 & \textbf{128}\\
64 & 384 & 336 & \textbf{304} & \textbf{304}\\
128 & 896 & 800 & 736 & \textbf{704}\\
256 & 2,048 & 1,792 & 1,664 & \textbf{1,600}\\
512 & 4,608 & 4,096 & 3,648 & \textbf{3,584}\\
1,024 & 10,240 & 8,960 & 8,320 & \textbf{7,936}\\
2,048 & 22,528 & 19,968 & 18,176 & \textbf{17,408}\\
4,096 & 49,152 & 43,008 & 38,912 & \textbf{37,888}\\
\bottomrule
\end{tabular}
\caption{Count~\cntG, exact $H_Nx$. ``Rate $7/8$ found in print'': the two formulas of~\cite{CFL}, page~969, which are
$\tfrac78N\lg N$ for even $k$ and $\tfrac78N\lg N+\tfrac18N$ for odd $k$ (one textbook layer and two
transforms of half the length). ``Layers of the $19$'': the $19$-operation circuit on three digits
at a time, then a textbook layer for one digit left over or the $7$-operation circuit for two. Bold: the
smallest of the row, when it is from this paper.}
\label{tab:G}
\end{table}

\paragraph{Where the gain comes from.}
A level of the textbook costs $N$. In the family a level costs $3N/4$: for each $d$, one hub and two pieces
in place of four pieces. The price is that a hub needs a partner, and the partner is a piece made two levels
below. So each piece is used twice: once to build the hub one level up, and once, two levels up, to turn a
hub into a piece. If one uses the $4$-point circuit on two digits at a time, with no sharing between
consecutive pairs of levels, the count obeys $T(N)=4T(N/4)+\tfrac74N$, which is the rate $7/8$
of~\cite{CFL}.


\section{What has been checked, and how}\label{sec:checked}

\paragraph{The statement checked in Lean.}
The Lean files are in the folder \verb|lean/| of the repository~\cite{Repo}; those written for this result
are in \verb|lean/Work/SmallWHT/|. The following statement was checked by Lean~4 (version 4.34.1) with the
library Mathlib:
\begin{flushleft}\ttfamily\small
theorem OAI.ExactFourier.WHT8Half.hadamard\_pow8\_circuit (j : $\mathbb{N}$) :\\
\hspace*{2em}$\exists$ C : Circuit (2 \textasciicircum{} (3 * j)),\\
\hspace*{3em}C.Computes (hadamard (3 * j)) $\wedge$ C.size $\le$ 20 * (j * 8 \textasciicircum{} (j - 1)) + 8 \textasciicircum{} j
\end{flushleft}
Here \verb|Circuit n|, \verb|size| and \verb|Computes| are OpenAI's
definitions~\cite{OAI1}: a circuit on $n$ inputs is a list of gates \verb|add i j|, \verb|sub i j| and
\verb|scale c i| with $c$ any complex number, together with a choice of $n$ outputs; \verb|size| is the number
of gates; \verb|C.Computes A| says that for every complex input $x$ the outputs are $Ax$. The model also offers the
constant $0$ as a value that costs nothing; this changes no count here. The subtraction
\verb|j - 1| is that of natural numbers (so $0-1$ is $0$). The prefix \verb|OAI| in the name is the name space
of OpenAI's Lean files, which my files extend; the files of this paper are mine, and this is not a project
of OpenAI. The matrix \verb|hadamard k| is defined in my files by
\begin{flushleft}\ttfamily\small
def hadamard (k : $\mathbb{N}$) : Matrix (Fin (2 \textasciicircum{} k)) (Fin (2 \textasciicircum{} k)) $\mathbb{C}$ :=\\
\hspace*{2em}fun j l => $\prod$ i $\in$ Finset.range k,\\
\hspace*{4em}if j.val.testBit i \&\& l.val.testBit i then (-1 : $\mathbb{C}$) else 1
\end{flushleft}
which is $H_{2^k}$ as defined in the Introduction. So the Lean statement is Theorem~\ref{thm:A}. A second
theorem in the same files states the bound as $\tfrac56N\lg N+N$ with real numbers.

\emph{What this covers.} In the Lean proof the $20$ operations on the right of Table~\ref{tab:h8}, the same
operations in another order, are data, and Lean checks that they compute $\tfrac12H_8$. The step to $8^j$ points uses OpenAI's Lean development of circuit costs
for tensor products; the $44$ files of theirs that the proof needs are in the repository, unchanged. The
folder \verb|lean/| of the repository was built once, on one Linux machine, in a new folder with a new
installation of Lean and with Mathlib's compiled files taken from its cache; the file \verb|VERIFY.md| of
the repository gives the commands. Lean reported no error, and the theorem depends only on the three standard
axioms \verb|propext|, \verb|Classical.choice| and \verb|Quot.sound|. A separate file holds the definitions
and the statement with no proof. A comparison script in the repository, written by the agents, checked that
the theorem that is proved is exactly that statement, with the same definitions. The comparator, which is
Lean's official program for this comparison, accepted the proof as a proof of that statement in an earlier
run, in the working copy in which the files were written; it was not run on the repository itself. The
comparator is meant to run inside a \emph{sandbox}, a
program that shuts the checked files off from the rest of the machine; that run used a simple stand-in for
it. Two controls were run on copies of files of the repository: with one gate of the $8$-point circuit changed,
Lean refused the proof, and with the bound $20$ changed to $19$, the comparison script refused the statement,
as the comparator had done in its run.

\emph{What it does not cover.} Nothing in count~\cntG{} is in Lean: not the $19$, not Theorem~\ref{thm:B}, not
Proposition~\ref{prop:lower}. Sizes that are not powers of $8$ (Corollary~\ref{cor:allk}) are not in Lean.
Above $8$ points the Lean proof builds the circuit inside the logic and writes out no list of gates. The
definition of \verb|hadamard| was written for this statement. A short Lean proof links it to a function
\verb|wht| that writes each output as a sum of the inputs with signs; both use the same formula for the
signs, so this is a check of the wording and not a second definition. No person has yet read the Lean files,
the statement included.

\paragraph{Exact replay by program.}
Two sets of Python programs check the circuits. Both work in exact arithmetic and run every circuit on all
$N$ unit inputs, or make one run that is equivalent to this, which determines a linear map.
\begin{itemize}\setlength{\itemsep}{1pt}\setlength{\parskip}{0pt}
\item The programs of the repository~\cite{Repo}, in the folder \verb|tools/circuits/|, write every circuit
as a text file, read it back and check it against the definition of $H_N$: the circuits of
Table~\ref{tab:h8}, the family, and the layered circuits of Section~\ref{sec:A}, for every size up to
$4{,}096$ points on one machine and up to $16{,}384$ on another. One command, run in the top folder of the
repository, repeats the run up to $4{,}096$ points: \verb|python3 -B tools/circuits/run.py|. The folder \verb|tools/circuits/fast/| has the circuits as working
programs in Python and C, with tests. The repository is under the Apache~2.0 licence.
\item A second program, \verb|check.py|, in the folder \verb|paper/| of the repository, was written
from the text of this paper. For Theorem~\ref{thm:B} it
follows the proof line by line and compares every line, when it is made, with the piece that the proof
says it is; it also confirms that the values are the same as those of the first generator. It was run for
every $k$ from $2$ to $14$. It also replays the layers of Corollary~\ref{cor:allk} for $k=3,\dots,14$,
recomputes every number in the tables, tests that every line of the $19$ and of the family is one operation
of~\cite{CFL} (Remark~\ref{rem:cfl}), and prints the numbers on the count of~\cite{GA} that the Introduction
quotes. It needs the Python library numpy. In the top folder of the repository the command is
\verb|python3 -B paper/check.py 12 tools/circuits|; the number is the largest~$k$.
\end{itemize}
A third program, \verb|lowerbound_check.c|, in the same folder \verb|paper/|, is about the proof of
Proposition~\ref{prop:lower}. Step~3 of that proof shows that the first $N-1$ lines of a circuit with $2N-2$
lines form a tree whose leaves are the inputs. The program goes through all such trees with $4$ and with $8$
leaves ($15$ and $135{,}135$ trees, each with every row at the root) and confirms Steps~4 and~5 at those
sizes: no case reaches more than one further row.

A replay is a check by a program, not a proof. For sizes above $16{,}384$ points the statements rest on
the proofs in this paper, which use nothing from outside it; Proposition~\ref{prop:19} can be checked by
hand from Table~\ref{tab:h8}.

\paragraph{The proofs.}
The proofs in this paper were written by the AI agents of Section~\ref{sec:found}. The proof of
Theorem~\ref{thm:B} was read by a second agent, and \verb|check.py| follows it line by line up to $16{,}384$
points. Proposition~\ref{prop:lower} and its proof were written on the morning of 11 October 2026, after the
literature search; a second agent read the proof, and the only check of it by program is that of
\verb|lowerbound_check.c|, for Steps~4 and~5 at $4$ and $8$ points.

\section{Remarks}\label{sec:remarks}

\subsection*{A lower bound, and the best counts I know at small sizes}

\begin{proposition}\label{prop:lower}
Let $N=2^k\ge4$. Every circuit of count~\cntG, with rational, real or complex constants, that has for each row
of $H_N$ a value equal to a non-zero multiple of that row applied to $x$, has at least $2N-1$ operations. The
same holds in count~\cntO, whose lines are special lines of count~\cntG.
\end{proposition}

\begin{proof}
Let the lines be $g_1,\dots,g_S$, in order. The \emph{support} of a value is the set of inputs with a non-zero
coefficient in it. A line \emph{depends on} its two operands and on everything they depend on; this is a matter
of the wiring only, so a constant equal to $0$ does not remove an operand. Call the \emph{ancestors} of a
line the lines that it depends on, itself included, and its \emph{cone} the set of inputs that it depends
on.
The support is contained in the cone.

\emph{Step 1.} A line with $t$ ancestors has a cone of at most $t+1$ inputs. Indeed, take as vertices the
$t$ ancestors and the inputs of the cone, and draw an edge from each operand of each ancestor to that
ancestor. This graph is connected, because every vertex leads to the line. It has at most $2t$ edges, and a
connected graph has at least one edge fewer than it has vertices. If the cone has exactly $t+1$ inputs, there
are exactly $2t$ edges and the graph is a tree.

\emph{Step 2.} Every output has all $N$ inputs in its support, so it is a line, and by Step~1 a line with full
support is $g_i$ with $i\ge N-1$. Different rows of $H_N$ are not multiples of each other, so the $N$ outputs
are $N$ different lines, all among $g_{N-1},\dots,g_S$. Hence $S\ge 2N-2$.

\emph{Step 3.} Suppose $S=2N-2$. Then the outputs are exactly $g_{N-1},\dots,g_{2N-2}$; the line $g_{N-1}$ has
full support, all of $g_1,\dots,g_{N-1}$ are its ancestors, and the graph of Step~1 is a tree with root
$g_{N-1}$. Call the inputs and $g_1,\dots,g_{N-2}$ the \emph{tree values}; the cone of an input is the input itself. In a tree the two operands of a line
have disjoint cones, and the cone of the line is their union. Let $g_{N-1}$ be $\lambda$ times row $a$,
written $\lambda\rho_a$. Going down from the root one sees that every tree value $w$ is a non-zero multiple
of $\rho_a$ restricted to the cone of $w$ (all other coefficients set to $0$): if a line $\alpha u+\beta v$ is
such a multiple, then on the cone of $u$ it equals $\alpha u$, because $v$ is zero there. The cones of the
tree values are \emph{nested}: any two are disjoint, or one contains the other. Each of them is neither
empty nor the set of all inputs.

\emph{Step 4.} Take an output $g_i=\alpha u+\beta v$ with $i\ge N$; it is $\lambda'\rho_b$ with $b\ne a$. Each
of $u,v$ is an earlier output or a tree value. They are not both outputs, because the rows of $H_N$ are
linearly independent. Two facts are used: two different rows of $H_N$ agree in exactly $N/2$ columns, and
no row is minus another, since all rows begin with $+1$.
If $u$ is an output, a multiple of $\rho_e$ with $e\ne b$, and $v$ is a tree value with cone $C$ (or the
other way round), then outside
$C$ the row $\rho_b$ is a multiple of $\rho_e$, so $\rho_b=\varepsilon\rho_e$ there with one sign
$\varepsilon$. On $C$ the difference $\rho_b-\varepsilon\rho_e$ is a multiple of $\rho_a$. It is not the zero
multiple, since $\rho_b\ne\varepsilon\rho_e$; so it is nowhere $0$ on $C$. Hence $C$ is exactly the set of the
$N/2$ columns where $\rho_b$ and $\varepsilon\rho_e$ differ, and on $C$ the difference is $2\rho_b$, so $\rho_b$
is a multiple of $\rho_a$ on $C$.
If $u$ and $v$ are both tree values, their cones cover all inputs, because $g_i$ has full support; so,
being nested, they are complementary
sets $C$ and $C'$; then $\rho_b=\pm\rho_a$ on $C$ and $\rho_b=\mp\rho_a$ on $C'$, and both sets have $N/2$
elements.
In both cases some tree value has a cone $C$ of size $N/2$ on which the product $\rho_a\rho_b$, taken column
by column, is constant. This product is a row of $H_N$ other than row $0$, so it has $N/2$ entries of each
sign, and it is therefore constant on $C$ and on the complement of $C$, with opposite signs.

\emph{Step 5.} Among nested sets at most two have size $N/2$, and if there are two they are complementary.
A set and its complement give the same product row, hence the same $b$. So at most one output besides
$g_{N-1}$ can be made. But there are $N-1\ge3$ of them. So $S\ge 2N-1$.
\end{proof}

The bound $2N-2$ is the obvious one: $N-1$ lines are needed before any value can contain all inputs, and
each further output needs a line. The proposition adds one. The bound is elementary, and no search of the
literature was made for it, so I do not say that it is new. The paper~\cite{CFL} claims no lower bound and
does not say whether $7$ is the least at $4$ points.

For $N=4$ the bound is $7$, so by Proposition~\ref{prop:lower} no circuit of count~\cntG{} for $4$ points has
fewer operations than the $7$-operation circuits. For $N=8$ and $N=16$ the bound
is $15$ and $31$, and the best circuits I know have $19$ (Section~\ref{sec:eight}) and $52$
(Theorem~\ref{thm:B}) operations.
The exact minimum at $8$ points is open: it is between $15$ and $19$. The searches described in
Section~\ref{sec:found} found no circuit with $18$ operations at $8$ points and none with $51$ at $16$, which
proves nothing. I found no exact minimum for $8$ or $16$ points in print.

\subsection*{The count of the 1994 paper}

\begin{remark}[every line is one multiply/add]\label{rem:cfl}
One operation of~\cite{CFL} is $\pm a\pm bc$, with any of the four choices of the two signs, or a plain
addition or multiplication. Every line of the count-\cntG{} circuits of this paper is $\pm u\pm v$ or
$\pm u\pm 2v$: it has two terms and at most one multiplication by a constant. This can be read off
Table~\ref{tab:h8}, the procedure $E(X,s)$ and the proof of Theorem~\ref{thm:B}, and \verb|check.py| tests
it for every line of the $19$ and of the family that it makes. Some lines have a minus sign on the term that
is not doubled: they are $-u+2v$ or $2u-v$. There are $2$ such lines among the $19$ (the lines $y_0$ and
$y_5$) and $12$ among the $52$ of the family at $16$ points. Such a line is the form $-a+bc$, so it is one
operation of~\cite{CFL} as well. Hence $19$, $52$ and $\tfrac34N\lg N+\tfrac14N$ are counts in the sense
of~\cite{CFL} as they stand. Count~\cntG{} as the Introduction defines it is wider in two ways: it allows a
constant on both operands, and it lets a multiplication on its own go free, where~\cite{CFL} charges one.
The count-\cntG{} circuits of this paper use neither.
\end{remark}

\subsection*{Two remarks without proofs}

\begin{remark}[other $2\times2$ matrices]
$H_N$ is the $k$-fold tensor power of the matrix with rows $(1,1)$ and $(1,-1)$: it is made from that
matrix by $k-1$ steps of the block rule for $H_{2N}$ in the Introduction. Put another $2\times2$ matrix
with four non-zero entries in its place; after scaling its rows and columns it has rows $(1,1)$ and
$(1,\tau)$. Programs of the agents report that, with $\tau$ kept as a symbol, the wiring of the family of
Section~\ref{sec:B} at $4$, $8$ and $16$ points, and the wiring of the $19$-operation circuit after a renaming
of the inputs and outputs and a scaling of each output, compute the corresponding transform with other
constants. This is computed, not proved, and no person has checked these programs. If it holds, the saving
is not special to the Walsh--Hadamard transform.
\end{remark}

\begin{remark}[a restricted class]
Every value of the circuits in this paper is, up to one factor, a vector with entries $0,\pm1$. Call the
number of an input, written with $k$ binary digits, its \emph{address}. In my working
notes there is an argument that if in addition the support of every value is a node of one fixed tree that
cuts the set of all addresses into two halves by the value of one digit, and each half again, down to single
addresses, then $N-1+\tfrac12N\lg N$ operations are needed: $7$, $19$ and $47$
at $4$, $8$ and $16$ points. The $19$-operation circuit is of that kind. No person has checked this
argument, so I state it as a direction and not as a result.
\end{remark}

\subsection*{The circuits as hardware}

\begin{remark}[count \cntG{} as a count of adders; measured, not proved]\label{rem:hardware}
In hardware built from \emph{adders} (an adder adds or subtracts two binary numbers) doubling is a shift by
one place, which is wiring. So each line of the count-\cntG{} circuits of this paper is one adder, and
their count is the number of adders: $19$ against the textbook's $24$ at $8$ points. The agents of
Section~\ref{sec:found} built the circuits and the textbook network for \emph{FPGAs}, chips whose logic is
set by loading a file into them and which are made of units called \emph{logic cells}. They used two
open-source tools: Yosys~0.69 for synthesis, which turns a design into logic cells, and nextpnr~0.11.1 for
placement and routing, which fits the cells onto the chip and wires them. Placed and routed for a Lattice
iCE40HX8K, as one block of adders with no \emph{pipeline registers} inside (rows of one-bit memories
between the levels of adders, which let the clock run faster), with $8$-bit inputs and one output
register, the $19$-adder circuit takes $274$ logic cells where the textbook network takes $349$, which is
$75$ fewer ($21.5\,\%$), or $273$ against $337$ ($64$ fewer, $19.0\,\%$) with the synthesis option
\texttt{-dff}. The
family of Theorem~\ref{thm:B}, built as such a block with signed $9$-bit inputs and no output register,
takes $17.5\,\%$ fewer cells than the textbook network at $16$ points, rising to $23.1\,\%$ fewer at $256$,
on the iCE40 (counted before placement; from $128$ points on the designs do not fit that chip), and
$19.1$ to $22.0\,\%$ fewer on a Lattice ECP5. With an output register and $8$-bit inputs the $16$-point
circuit takes $877$ cells against $969$ ($9.5\,\%$ fewer), and the $56$-adder design of~\cite{CFL} takes
$720$.

The saving has a price. The circuits are twice as deep ($6$ levels of adders against $3$ at $8$ points),
so the tools estimate the highest clock rate of the $8$-point block, the number of results it can take a
second, at about $72$\,MHz against about $97$\,MHz for the textbook (timing estimates, not measurements),
and the family's blocks as about $30\,\%$ slower
than the textbook's on the iCE40 and about $37\,\%$ slower on the ECP5. Once pipeline registers are added
for speed, the textbook is the smaller design: above about $85$\,MHz at $8$ points (with a register after
every adder, $677$ cells against $363$), and with a register after every level of adders the family needs
$67$ to $70\,\%$ more iCE40 cells than the textbook. The $19$-adder circuit is also not the smallest
$8$-point design on this chip: the $22$-adder step of Alman and Rao~\cite{AR} takes $244$ cells under the
same conditions. With these tools a subtraction costs more than an addition on this chip, and about half
of the saving over the textbook comes from having $9$ subtractions against $12$; their step has one.
Without the output register the family takes $10$ to $14\,\%$ fewer cells than the design of~\cite{CFL} on
the ECP5 and about as many on the iCE40 (between $2\,\%$ fewer and $10\,\%$ more, depending on the size).
% The next sentence is kept on one line of this source.
On an iCE40-HX8K FPGA board running at 12 MHz, the 19-adder 8-point Walsh-Hadamard circuit was run side by side with the textbook circuit on 597{,}750{,}000 test inputs, and the two gave the same outputs every time; a deliberately broken copy of the circuit, run alongside, was caught, which shows the test can fail. This checks only that the circuit computes the right values on a real chip: it says nothing about size or speed, and in this pipelined build the 19-adder circuit is the larger of the two designs. The lines that the board sent are in the repository, in \texttt{tools/circuits/fpga/board/run-2026-10-11/}.
These are measurements and estimates by the agents with one open-source tool chain; the chip makers' own
tools were not run, and nothing was measured for custom silicon or for power. The files and the conditions
are in the folder \verb|tools/circuits/fpga/| of the repository~\cite{Repo}; those of the timings of
Section~\ref{sec:not} are in \verb|tools/circuits/fast/wht56/results/ffht/|.
\end{remark}


\section{How the circuits were found, and the use of AI}\label{sec:found}

The circuits were found on the night of 10 to 11 October 2026 by computer searches that were written and
run by AI coding agents (Anthropic's Claude) under my direction. They were a by-product of a search for
something else, a transform of linear size, which found nothing of that kind.
The working reports of that search, about 180 of them with their programs, are in the folder
\texttt{waves/} of the repository~\cite{Repo}; its page \texttt{WHERE-THE-CIRCUITS-CAME-FROM.md} names
the report behind each step described here, and \texttt{LEADS.md} lists what the search left open.
One step of that search was
to try to prove that $N\lg N$ operations are needed when a circuit holds only pieces in the sense of
Section~\ref{sec:B}. A local search in a slightly larger set of values, started from a $22$-operation
circuit, fell to $20$ operations at $8$ points. The $20$ contained a value that is not a piece: the hub of
Lemma~\ref{lem:hub}. Used as a rule at every size, it is the family of Theorem~\ref{thm:B}. A random search
over sets of helper values (simulated annealing), again started from the $22$, then fell to $19$ in each of
four runs and never to $18$. The Lean proof of Theorem~\ref{thm:A}, the programs, the literature search,
the proofs and this text were also written by the agents. I am responsible for the
content of this paper.
The agents also ran the programs and the checks of the repository, and made the measurements of
Section~\ref{sec:not} and Remark~\ref{rem:hardware}, at my direction, on a Mac mini of mine
and on a Linux machine that I rented. No person has yet read the Lean files or checked the proofs line by line.

\emph{The literature search} was also made by the agents: the papers that the citation indexes list as citing
\cite{AR} or \cite{CFL}, most of them seen by title and abstract only; the text of \cite{AR}, \cite{Feig},
\cite{VP} and \cite{GB}; several dozen queries to arXiv and to web search engines; and seven patents, read
through a tool that summarises pages. The two older papers \cite{CFL} and \cite{GA} were read by the
agents, from copies I obtained. The first has two pages, and both were read. Of the second the agents read
page~172 (the introduction) and pages~178 to~183: the comparison tables, the paragraph on the count for the
Hadamard transform, the conclusions, and the appendix with the proof of the count and the programs.
Pages~173 to~177 were not read. By the outline on page~172 they hold the part on the Fourier transform and
also the section that sets out the algorithm for the Hadamard transform, so the algorithm is taken here from
the proof and the program in the appendix. Nothing
found has $8$ points below $22$ operations, $16$ points below $56$, a rate below $7/8$ in count~\cntG, or a rate
below $23/24$ in count~\cntO; the count stated in~\cite{GA} is not such a result, for the reason given in the
Introduction.

\emph{What limits the word ``new''.} The paper
\cite{Samadi} is known from its abstract only. The same holds for three works of Sarukhanyan, Agaian, Astola and
Egiazarian on Hadamard transforms and multiply-add architectures~\cite{SAAE03,SAAE05,ASEA}; their abstract
speaks of beating the classical transform for orders $96$ and above, and what is compared could not be told
from it.
A survey by I.~S.~Sergeev~\cite{Sergeev} has a section called ``Linear circuits for Sylvester matrices''
($H_N$ is also called the Sylvester matrix); only the first $49$ pages of the survey were read, and that
section is not in them. The comment of Frigo and Johnson~\cite{FJ} is known from its abstract only, and
the paper of Morgenstern~\cite{Morgenstern} from later accounts. Five pages of~\cite{GA} were not read, as
said above. No search was made for lower bounds like Proposition~\ref{prop:lower}. Literature in
languages other than English was seen only as far as the indexes returned it. If you
know a smaller count in print, I would be glad to hear of it and will credit it.

\section{What this is not}\label{sec:not}

\begin{itemize}\setlength{\itemsep}{1pt}\setlength{\parskip}{0pt}
\item \textbf{It is not a better exponent.} Every count here is a constant times $N\lg N$. The results
of~\cite{OAI1,OAI2,RepoPS} are of smaller order. The count-\cntG{} circuits of this paper cannot reach a smaller
order as they are: their constants are $\pm1$ and $\pm2$, and for circuits whose constants are bounded,
Morgenstern's determinant argument~\cite{Morgenstern} gives a lower bound of order $N\log N$.
\item \textbf{It is not a statement about running time.} No result of this paper is a timing, and where
time was measured, fewer operations did not mean less time. The test that counts is against FFHT, a
well-known fast open-source library for this transform, whose code is written out as \emph{vector
instructions} (processor instructions that work on several numbers at once). The $8$-point step was
written into the code that FFHT's own generator produces, in place of its block of $24$ additions and
subtractions, and nothing else was changed. On the one machine that was timed, this did not make FFHT
faster at any size from $2^6$ to $2^{26}$ points in single precision or from $2^5$ to $2^{26}$ points in
double precision: set against a second build of FFHT's unchanged code, all $129$ combinations of size and
form took more time.\footnote{Conditions. One rented virtual machine with the AVX2 and FMA sets of vector
instructions, which reports itself as an Intel Xeon at $2.20$\,GHz, with other programs running on it at
full load the whole time; the compiler gcc~13.3.0, for the plain C with \texttt{-O3 -march=native}; FFHT,
\url{https://github.com/FALCONN-LIB/FFHT}, at its last version (March 2019), built by its own Makefile,
with the plan of passes it ships; one thread, kept on one processor; wall-clock time, the best of $15$
repeats of a loop of $0.2$\,s or more ($9$ repeats where a size was timed a second time, and for the plain
C); 11 October 2026. The three forms: the $20$ operations on the right of Table~\ref{tab:h8} with the
output left divided by a power of two; the same with multiplications added so that the output is exact;
the $19$ lines on the left with fused multiply-adds, exact. From $2^7$ points ($2^6$ in double precision)
to $2^{19}$ points the changed FFHT took between $1$ and $24\,\%$ more time, depending on the form and the
precision, and from $2^{20}$ points $2$ to $7\,\%$ more on average, which for the fused form is too small
a difference for one run to establish. Plain C: $7$ to $37\,\%$ less time with exact output at the sizes
named, $1$ to $13\,\%$ more at every other size from $2^4$ to $2^{26}$ in double precision, and $2$ to
$51\,\%$ more at every size from $2^3$ to $2^{27}$ in single precision; $20$ operations in place of $24$
can save $17\,\%$ at most. The timings were made by the agents of Section~\ref{sec:found}, in one run unless
said otherwise, and have not been repeated on a second machine, with a second compiler, or by anyone
else.}
One of the three forms that were tried has the $19$ lines of count~\cntG{} as $19$ instructions: each line
with a doubling is done by a \emph{fused multiply-add}, a single instruction that computes $a+b\times c$.
It was not faster either. The likely reason is the order in which the operations depend on each other: in
the textbook block no result is more than $3$ operations away from the loaded numbers, and in the new step
the last results are $7$ operations away, $6$ in the fused form. What happens inside the processor was not
measured, so this is an explanation that fits, not a finding.
In my own plain C, where the compiler writes the vector instructions, the new step took less time than the
textbook step on that machine only in double precision at the sizes $2^6,2^9,\dots,2^{27}$; it took more
time at every other size from $2^4$ to $2^{26}$, and in single precision at every size. At $2^6$ to
$2^{18}$ points that gain is larger than replacing $24$ operations by $20$ can give, so it comes from how the compiler treats one pass
of my code and not from the operation count; the cause has not been found. Other timings of plain C, on
that machine and on a Mac mini, are beside the programs in the repository~\cite{Repo}; in no build did the
time follow the operation count. Only the $8$-point step was timed, not the larger circuits of the family.
The 1994 paper shows the same gap between counting and timing: its count
predicts a gain of $12.5\,\%$ over an ordinary routine of the same kind, and the authors measured $5.6\,\%$
to $7.4\,\%$ on their machine, an IBM RISC System/6000 Model~530, at $64$ to $4{,}096$
points~\cite[pages~969 and~970]{CFL}.
\item \textbf{It is not known to be optimal.} The only lower bound proved here is the $2N-1$ of
Proposition~\ref{prop:lower}, one more than the obvious $2N-2$. In count~\cntO{} I do not know whether the exact transform on fewer than $512$ points can be computed with fewer
than $N\lg N$ operations. Nor do I know whether the rate $5/6$ can be lowered. An attempt to bring the rate
$3/4$ of the family into count~\cntO{} found nothing below $5/6$ (Remark~\ref{rem:factor}: computations and an
argument of the agents, not in Lean). A direct search for a $16$-point circuit with $53$ or fewer operations
of count~\cntO{} has not been made.
\item \textbf{The method is not mine.} Splitting a small transform into a part of low rank and a sparse
part, and the single halving, are how Alman and Rao obtain their $8$-point step. The $4$-point circuit
with $7$ operations is due to Coppersmith, Feig and Linzer~\cite{CFL}. What the searches added is which
values to share.
\end{itemize}

\begin{thebibliography}{99}
\setlength{\itemsep}{1pt}\setlength{\parskip}{0pt}

\bibitem{ASEA}
S.~Agaian, H.~Sarukhanyan, K.~Egiazarian and J.~Astola, \emph{Hadamard Transforms}, SPIE Press, Bellingham,
2011. (Abstract only.)

\bibitem{Alman}
J.~Alman, Kronecker products, low-depth circuits, and matrix rigidity, in \emph{Proceedings of the 53rd Annual
ACM SIGACT Symposium on Theory of Computing (STOC 2021)}, pages 772--785.
\href{https://doi.org/10.1145/3406325.3451008}{doi:10.1145/3406325.3451008}; arXiv:2102.11992.

\bibitem{AR}
J.~Alman and K.~Rao, Faster Walsh-Hadamard and discrete Fourier transforms from matrix non-rigidity, in
\emph{Proceedings of the 55th Annual ACM Symposium on Theory of Computing (STOC 2023)}, pages 455--462.
\href{https://doi.org/10.1145/3564246.3585188}{doi:10.1145/3564246.3585188}; arXiv:2211.06459, version~2.

\bibitem{CFL}
D.~Coppersmith, E.~Feig and E.~Linzer, Hadamard transforms on multiply/add architectures, \emph{IEEE
Transactions on Signal Processing} 42(4) (1994), 969--970.
\href{https://doi.org/10.1109/78.285664}{doi:10.1109/78.285664}.

\bibitem{Feig}
E.~Feig, Complexity theory of transforms in signal processing, Chapter~9 of \emph{The Digital Signal
Processing Handbook}, edited by V.~K.~Madisetti and D.~B.~Williams, CRC Press, Boca Raton, 1999.

\bibitem{FJ}
M.~Frigo and S.~G.~Johnson, Comments on ``Split manageable efficient algorithm for Fourier and Hadamard
transforms'', \emph{IEEE Transactions on Signal Processing} 55(4) (2007), 1547--1548.
\href{https://doi.org/10.1109/TSP.2006.889087}{doi:10.1109/TSP.2006.889087}. (Abstract only.)

\bibitem{GA}
A.~M.~Grigoryan and S.~S.~Agaian, Split manageable efficient algorithm for Fourier and Hadamard transforms,
\emph{IEEE Transactions on Signal Processing} 48(1) (2000), 172--183.
\href{https://doi.org/10.1109/78.815487}{doi:10.1109/78.815487}. (Pages 172 and 178 to 183 only.)

\bibitem{GB}
A.~M.~Grigoryan and V.~S.~Bhamidipati, Method of flow graph simplification for the 16-point discrete Fourier
transform, \emph{IEEE Transactions on Signal Processing} 53(1) (2005), 384--389.
\href{https://doi.org/10.1109/TSP.2004.838970}{doi:10.1109/TSP.2004.838970}.

\bibitem{Morgenstern}
J.~Morgenstern, Note on a lower bound on the linear complexity of the fast Fourier transform, \emph{Journal
of the ACM} 20(2) (1973), 305--306.
\href{https://doi.org/10.1145/321752.321761}{doi:10.1145/321752.321761}. (Not read; known through
later accounts.)

\bibitem{OAI1}
OpenAI, Finite tensor savings and exact Fourier circuits, preprint of 25 September 2026, with a Lean
development, in the folder \texttt{preprints/} of the repository \url{https://github.com/openai/math}.

\bibitem{OAI2}
OpenAI, An explicit power saving for the exact discrete Fourier transform, preprint of 25 September 2026, in
the folder \texttt{preprints/} of the repository \url{https://github.com/openai/math}.

\bibitem{Patent}
E.~Reshef and I.~Alrod, Reduced complexity fast Hadamard transform, United States patent 7,003,536~B2, filed
15 August 2002, granted 21 February 2006. (Read through a tool that summarises pages.)

\bibitem{Samadi}
S.~Samadi, Y.~Suzukake and H.~Iwakura, On automatic derivation of fast Hadamard transform using genetic
programming, in \emph{Proceedings of the 1998 IEEE Asia-Pacific Conference on Circuits and Systems
(APCCAS 1998)}, pages 327--330.
\href{https://doi.org/10.1109/APCCAS.1998.743763}{doi:10.1109/APCCAS.1998.743763}. (Abstract only.)

\bibitem{SAAE03}
H.~Sarukhanyan, S.~S.~Agaian, J.~T.~Astola and K.~O.~Egiazarian, Binary matrices, decomposition and
multiply-add architectures, in \emph{Image Processing: Algorithms and Systems II}, Proceedings of SPIE 5014
(2003), 111.
\href{https://doi.org/10.1117/12.473134}{doi:10.1117/12.473134}. (Abstract only.)

\bibitem{SAAE05}
H.~Sarukhanyan, S.~Agaian, J.~Astola and K.~Egiazarian, Decomposition of binary matrices and fast Hadamard
transforms, \emph{Circuits, Systems and Signal Processing} 24(4) (2005), 385--400.
\href{https://doi.org/10.1007/s00034-004-0811-y}{doi:10.1007/s00034-004-0811-y}. (Abstract only.)

\bibitem{Sergeev}
I.~S.~Sergeev, Etudes on the methods of designing efficient circuits, a survey of about 170 pages, version
of 28 May 2025, on ResearchGate. (Pages 1 to 49 only.)

\bibitem{Repo}
J.~Sussman, \emph{\reponame}, the repository of this paper, with the circuits, the programs, the Lean files
and the source of the paper, 2026. \url{\repourl}.

\bibitem{RepoPS}
J.~Sussman, \emph{wht-power-saving-lean}, an earlier repository of Lean proofs and programs, 2026.
\url{https://github.com/jacobalansussman/wht-power-saving-lean}.

\bibitem{VP}
Y.~Voronenko and M.~P\"uschel, Mechanical derivation of fused multiply--add algorithms for linear transforms,
\emph{IEEE Transactions on Signal Processing} 55(9) (2007), 4458--4473.
\href{https://doi.org/10.1109/TSP.2007.896116}{doi:10.1109/TSP.2007.896116}.

\end{thebibliography}

\end{document}
